\section{The Natural Logarithm-tmp}
\label{sec:natural-logarithm-1}
\markright{\sc{}The Natural Logarithm}
\subsection*{\underline{Introduction: John Napier}}
\begin{wrapfigure}[]{r}{2in}
  \captionsetup{labelformat=empty}
  \centerline{\includegraphics*[height=2.5in,width=2in]{../Figures/Napier.png}}
  \label{fig:Napier}
  \caption{\textcolor{blue}{John Napier, $1555-1617$}}
\end{wrapfigure}


In the $16$th and $17$th centuries the astronomical calculations
needed for navigation were growing increasingly complex. These
computations were done by hand, frequently on the deck of a
wind-tossed sailing ship, and often by people who did not fully
understand the underlying ideas. It was important that the
computations be as easy as they could possibly be, and still be
accurate.  The Scottish mathematician John Napier ($1550-1617$)
invented logarithms specifically to address this problem. Logarithms
have the remarkable property that by using them you can convert
multiplication to addition, division to subtraction, exponentiation to
multiplication and root-taking to division. In each instance
decreasing the complexity of the operation.

Napier described their properties in \foreign{Mirifici logarithmorum
  canonis descriptio}, published in $1614$, well before Newton or
Leibniz were born so logarithms actually predate the invention of
Calculus by several decades.  He coined the term ``logarithm'' from
the Greek words logos meaning ``reasoning,'' or `` reckoning'' and
arithmos meaning ``number.''  So to Napier logarithms were
``reckoning numbers'' which seems an apt description, given their
original purpose. Nowadays they are much, much more.


The ability to turn multiplication into addition and exponentiation to
multiplication was a major innovation because using simpler operations
meant fewer errors were made and the better chance a ship had of
actually reaching its destination.

We will discuss Napier and his work a little later. For now we focus
on the basic properties of  the natural logarithm that you will need
to solve problems in this course and after.
\begin{center}
  \begin{minipage}[h]{.6\linewidth}
    \centerline{\bf{}\underline{Properties of the Natural Logarithm}}
    \begin{enumerate}
    \item $\ln(e)=1$,\\ where $e$ is the base of the natural exponential,
      $e^x$.      
    \item $\ln(ab)=\ln(a)+\ln(b)$
    \item $\ln\left(\frac{a}{b}\right)=\ln(a)-\ln(b)$
    \item $\ln\left(a^b\right)=b\ln(a)$
    \end{enumerate}
  \end{minipage}
\end{center}
Like the differentiation rules you'll need to be very comfortable
using these properties.  The simplest use we can put a \term{natural
  logarithm} to is to aid in  solving an equation like:
$$
3^x=17.
$$
The difficulty here is how do we get the variable $x$ out of the
exponent. This is exactly what logarithms were invented to do. Notice
that by Property 4, $\ln(3^x)=x\ln(3)$, so we can solve our equation
like this:
\begin{align*}
  3^x&=17\\
  \ln(3^x)&=\ln(17)\\
  x\ln(3)&= \ln(17)\\
  x&=\frac{\ln(17)}{\ln(3)}.
\end{align*}
At this point we can use technology to find that
$\ln(17) \approx 2.8332$ and $\ln(3)\approx1.0986$ so that
$\frac{\ln(17)}{\ln(3)}\approx2.5789$.

Of course if that's all we need logarithms for we could just stop
here. With practice anyone can go through the motions and solve
equations without actually understanding what they are doing just like
the differentiation rules. As we said, that was the whole point of
inventing logarithms in the first place, but using them is just the
first step in understanding them and our true goal is understanding.  

\begin{ProficiencyDrillSection}
  \begin{ProficiencyDrill}
    Solve each of the following equations:
    \begin{enumerate}[label={\bf{}(\alph*)}]
      \begin{multicols}{2}
      \item $2^{x}=3^x$
        \solution{$x=0$}
      \item $e^{x-5}=e^5$
        \solution{$x=10$}
      \item $7^{x^2-x}=e$
      \end{multicols}
    \end{enumerate}
  \end{ProficiencyDrill}
\end{ProficiencyDrillSection}

\subsection*{\underline{Inverting the Natural Exponential}}
\label{sec:underl-logar-as}
The natural logarithm is the final ``elementary'' function we will
consider.  It is most simply defined as the inverse of the natural
exponential, but that is not a very helpful definition.  We will need
the natural logarithm for much the same reasons that we needed the
inverse trig functions in Trigonometry so a quick refresher on those
will be helpful.


Recall that if we have a particular angle, $\theta,$ it is easy enough
to understand what $\sin(\theta),$ $\cos(\theta),$ $\tan(\theta),$
etc. mean and to compute each of these for a particular value of
$\theta$. But sometimes this is not enough.

Sometimes we have the value of sine (or cosine, or tangent, $\ldots$)
of an angle $\sin(\theta),$ but we need to find $\theta$ itself. For
example, if we were trying to solve the equation
$$
4\sin(2x+1)=3
$$
we'd first reduce this to
$$
\sin(2x+1)=3/4.
$$
But if we don't know how to ``undo'' the sine we're stuck. This is
what  \term{inverse sine} function was invented for. It is
the function that ``undoes'' the sine. To solve our equation we simply
apply $\inverse\sin()$ to both sides. This gives:
$$
\inverse\sin(\sin(2x+1))=\inverse\sin(3/4).
$$
\emph{By definition} the \term{inverse sine} ``undoes\aside{You
should read the symbol $\inverse\sin(x)$ as ``the angle whose sine is
$x$.''
Reading the formula $\sin\left(\inverse\sin(x)\right)=x,$ aloud in this way
gives us ``the sine of the angle whose sine is $x.$''
Similarly $\inverse\sin(\sin(x))=x$ is read as ``the angle whose sine
is the sine of $x$ is $x$.''}'' the sine. 

The basic idea of \alert{function inversion} is really as simple as
that. If $f(x)$ is the name of a function, then $\inverse{f}(x)$ is
the name of the \emph{different but related} function which ``undoes''
$f.$ This is why, for example,
\begin{align*}
  \inverse\sin(\sin(x))&=x \\
  \inverse\cos(\cos(x))&=x \\
  \inverse\tan(\tan(x))&=x, \text{ etc. }
\end{align*}

Ok. Sure. We all know that using the inverse trigonometrc functions is
a bit harder than this. But the complications you remember from
Trigonometry didn't come from function inversion. The complications
came from the formalities of the underlying, abstract notion of a
\term{function.}  Fortunately, when we invert the natural exponential
none of those formalities come into play, so we needn't bother with
them. For now.

The \term{natural logarithm} function is, by definition, the inverse
of the \term{natural exponential} in precisely the same way that
$\inverse\sin(x)$ is, by definition, the inverse of $\sin(x)$ .
Unfortunately the notation that is used for the \term{natural
  exponential} and the \term{natural logarithm} is not helpful.

\begin{center}
  \begin{minipage}[h]{.85\linewidth}
    \begin{mynotation}{Problems with Inverse Function}
      It has become standard practice to name functions using the
      first three letters of the name of the function. That is why the
      trigonometric functions are $\sin()$, $\cos()$, $\tan()$,
      $\cot()$, $\sec()$, and $\csc()$. By that standard we should
      \emph{properly} denote the \term{natural exponential} with
      $\exp(x)$. We actually will use this notation when we want to
      emphasize the functional nature of the \term{natural
        exponential}, but most of the time it is more helpful to think
      of the \term{natural exponential} as ``$e$ raised to a power''
      and to denote it as $e^x$ which, unfortunately, doesn't
      \emph{look} like a function at all.      

      It is also fairly standard practice to denote the inverse of a
      given function with the exponent $-1$. This is how the inverse
      trig functions are often denoted: $\inverse\sin()$,
      $\inverse\cos()$, $\inverse\tan()$, $\inverse\cot()$,
      $\inverse\sec()$, and $\inverse\csc()$.  By that standard we
      should \emph{properly} denote the the \term{natural logarithm}
      by $\inverse\exp()$, but this is almost never done.

      Since the logarithm has its own name (unlike the inverse trig
      functions) it would be reasonable to denote it with the first
      three letters of the word ``logarithm'': $\log(x)$. Most
      mathematicians actually will use this symbol when we talk
      amongst ourselves, because for us the \term{natural logarithm }
      is the only logarithm function that really matters.

      However, since the time of Napier the symbol $\log(x)$ has
      always denoted what are called \term{common logarithms} (we'll
      have more to say about these shortly). As a result the symbol
      that is most frequently used by scientists and engineers (and by
      mathematicians when we are talking to scientists and engineers.)
      to denote the \term{natural logarithm} is $$\ln(x).$$

      None of this affects what the \term{natural logarithm} is, how
      we use it, or how we think about it. It only affects how we
      write it. The \term{natural logarithm} ($\ln(x)$) is the inverse of the \term{natural
        exponential} ($e^x$) and you should think of it as such.

      However, rather than writing
      $$
      \inverse\exp(\exp(x))=x, \text{  or  } \exp(\inverse\exp(x))=x
      $$
      as we would with almost any other function/inverse pair for the
      \term{natural exponential}/\term{natural logarithm} pair we
      write instead
      $$
      \ln(e^x)=x, \text{  and  } e^{\ln(x)}=x
      $$
      because that is the notation that everyone uses and
      understands. When using this notation it is often helpful to
      think of $\ln(x)$ as the function that peels the exponent off of
      $e^x$. This is actually what gives us Property 1:
      $\ln(e)=\ln(e^1)=1$.

\end{mynotation}
\end{minipage}
\end{center}

\begin{ProficiencyDrill}
  Evaluate each of the following:
  \begin{enumerate}[label={\bf{}(\alph*)}]
    \begin{multicols}{2}
    \item $\ln(e^2)$ \solution{$2$}{}
    \item $\ln(e^{\sin(t)}$ \solution{$\sin(t)$}{}
    \item $e^{\ln(2)}$ \solution{$2$}
    \item $e^{\ln(x+y)}$ \solution{$x+y$}
    \end{multicols}
  \end{enumerate}
\end{ProficiencyDrill}


\subsection*{The Derivative of the Natural Logarithm}
We don't normally  think of it this way but  the natural exponential
changes addition to multiplication and multiplication to
exponentiation in the following manner:
\begin{equation}
e^{\overbrace{a+b}^{\text{addition}}}=\underbrace{e^ae^b}_{\text{multiplication}},\text{ and } \underbrace{e^{ab}}_{\text{multiplication}}=\underbrace{(e^a)^b}_{\text{exponentiation}}\label{eq:ExpDefProp}
\end{equation}
or in function notation
$$
\exp(\underbrace{a+b}_{\text{addition}})=\underbrace{\exp(a)\exp(b)}_{\text{multiplication}},\text{
  and
} \exp(\hskip-.25in\underbrace{ab}_{multiplication}\hskip-.25in)\hskip.25in=\underbrace{(\exp(a))^b}_{\text{exponentiation}}.
$$

Since the natural logarithm is the inverse function of the natural
exponential it ``undoes'' everything that the natural exponential
``does.'' That is, it will change multiplication to addition and
exponentiation to multiplication as follows:
\begin{equation}
  \underbrace{\ln(ab)}_{\text{multiplication}}=\underbrace{\ln(a)+\ln(b)}_{\text{addition}} \text{\hskip5mm{}and }
  \underbrace{\ln(a^b)}_{\text{exponentiation}}=\underbrace{b\ln(a)}_{\text{multiplication}}.
  \label{eq:LogProps}
\end{equation}

\begin{embeddedproblem}
The properties in equation~\ref{eq:LogProps} are two of the properties
listed in the table of properties at the beginning of this section.
Property 3 in that table is given just for convenience. It is really
an immediate consequence of Property 2 and Property 4.

Derive Property 3 from the others. \hint{Find a way to express
  $\frac{a}{b}$ as a product.}
\end{embeddedproblem}


Since the \term{natural logarithm} is the inverse of the \term{natural
  exponential} we can find its derivative by using the same trick that
we used to find the derivatives of the inverse trig functions. We
start with the identity
\begin{align*}
  e^{\ln(x)}&=x\\
   \intertext{and  make this easier on our eyes with the substitution
   $y=\ln(x).$ Differentiating this, we have}
   e^y\dx{y} &= \dx{x}\\
   \dx{y}&=\frac{1}{e^y}\dx{x}.\\
   \intertext{Making the  substitution $y=\ln(x)$ once more we have}
   \dx{y}&=\frac{1}{e^{\ln(x)}}\dx{x},\\
   \intertext{ and since $e^{\ln(x)}=x$ we have}
   \dx{y}&=\frac{1}{x}\dx{x},\\
\end{align*}
and this is the last of the differentiation rules you need to
memorize. That was so easy it's kind of a letdown\footnote{}, isn't it?  This is the last
differentiation rule you need to
memorize:
$$
\dx{(\ln(x))}=\frac{1}{x}
\dx{x}=x^{-1} \dx{x}.
$$


 Like all differentiation rules we need to use this one in tandem with
 all of the others so practice is required.

\subsubsection*{\underline{Two  Curious Facts}}
This differentiation rule is actually quite remarkable
for a couple of reasons.
\begin{description}
\item[Fact 1:] Recall that in
  section~\ref{sec:inverse-trig-funct} we pointed out that
  differentiating a function often results in another function of the
  same type.  The inverse trig functions do not have this property and
  now we see that the \term{natural logarithm} also breaks this
  pattern.

  Do you suppose there is any significance to this?
\item[Fact 2:] 
  Somewhat more remarkable is what happens when we try to run the
  differentiation process backwards.

  Observe that if $y=x^\alpha,$ then
  $\dfdx{y}{x}=\alpha x^{\alpha-1}.$ This is, of course, just the
  Power Rule. In words, it says that the derivative of any monomial
  $(x^\alpha)$ is another monomial. But what happens if we try to go
  the other direction?

  Consider: If $\dfdx{y}{x}=x^2$ the clearly $y=\frac{x^3}{3},$ and in
  general if $\dfdx{y}{x}=x^\alpha$ then
  $y= \frac{x^{\alpha+1}}{\alpha+1}.$ This is true as long as
  $\alpha\ne -1.$ In that case running the Power Rule backwards would
  give us $y= \frac{x^{-1+1}}{-1+1}$ which won't work because in the
  denominator we have $-1+1=0.$

  To put all of this in words, if we start out with a monomial
  derivative and ask for the original function we find that it will
  also be a monomial for every possible value of the exponent
  \emph{except} $-1.$ In that case the original function would be
  $y=\ln(x).$ It seems very strange that there would be this one,
  single exception doesn't it?

  Or, to come at this in a different way maybe this says that the
  logarithm function is somehow related to monomials? What do you
  think?
\end{description}

 \begin{ProblemSection}

\begin{myproblem}
  Compute each of the following:
  \begin{enumerate}[label={\bf{}(\alph*)}]
    \begin{multicols}{3}
    \item $\dx{\left(\sqrt{x^2+1}\cdot\ln\sqrt{x^2+1}\right)}$
    \item $\dx{\left(\dfrac{\ln(x)}{e^x}\right)}$
    \item $\dx{\left(\tan(ln(x))\right)}$
    \item $\dx{\left(\ln(\tan(x))\right)}$
    \item $\dx{\left(\inverse\tan(\ln(x))\right)}$
    \item $\dx{\left(x\ln(x)-x\right)}$
    \item $\dx{\left(\ln(\ln(\sec(x)))\right)}$
    \item $\dx{\left(\ln(\sin(x^2+1))\right)}$
    \item $\dx{\left(\sec(x)\tan(x)\right)}$
    \item $\dx{\left(\frac{e^{x^2}\sin(x)}{\sqrt{x^3-7x+4}}\right)}$
    \item
      $\dx{\left(\left(\frac{(2x+3)(4x+5)^{20}}{x^{\frac12}-\ln(x)}\right)^{\frac43}\right)}$
    \end{multicols}
  \end{enumerate}
\end{myproblem}


\begin{myproblem}{}
  Solve the following equations for $x.$
  \begin{enumerate}[label={\bf{}(\alph*)}]{}
    \begin{multicols}{3}
    \item $e^x=\frac{1}{\sqrt{e}}$
    \item $e^x=5$
    \item $7e^{x+3}=2$
    \item $\ln(2x-4)=0$
    \item $\ln(2x)-\ln(x+1)=4$
    \item $2\ln(x)=\ln(3x-2)$
    \end{multicols}
  \end{enumerate}
\end{myproblem}

\begin{myproblem}
  Solve each of the following for $c:$
  \begin{enumerate}[label={\bf{}(\alph*)}]
    \begin{multicols}{3}
    \item $-\ln(2c)=8$
    \item $-\ln\left(\frac{1}{2c}\right)=8$
    \item $\ln(2c)+\ln(2)=8$
    \item $\ln(2x)\ln(x)=\ln(c)$
    \item $\ln(5x)=\ln(x)+c$
    \end{multicols}
  \end{enumerate}
\end{myproblem}

\begin{myproblem}
  Solve for $x$:
  \begin{enumerate}[label={\bf{}(\alph*)}]
    \begin{multicols}{2}
    \item $e^{2x}-3e^x+2=0$
    \item $(\ln(x))^2+3\ln(x)+2=0$
    \item $\sin(e^x)=0$
    \item $\cos(e^x)=0$
    \item $e^{2x}- e^x-2=0$
    \item $(\ln(2x)^2- \ln(x)-2=0$
    \item $(\ln(x))(\ln(2x))=1$
    \end{multicols}
  \end{enumerate}
\end{myproblem}

\begin{myproblem}{}
  For each of the following find $\d{y},$  and $\dfdx{y}{x}$.
  \begin{enumerate}[label={\bf{}(\alph*)}]{}
  \item $y=\ln(2x+2)$
  \item $y=\ln
    \left(
      \frac{(2x+1)^2(3x-2)}{x^2+1}
    \right)$ \hint{Think about properties of logarithms before you
      do this.} 
  \item $\ln(y)=\sin(e^x)$
  \item $\ln(x+y)=\inverse\tan(z)$
  \end{enumerate}
\end{myproblem}

\begin{myproblem}
  Show that if $y=(x+1)(x+2)(x+3)\ldots(x+n),$ then for $x\neq1, 2,
  3, \ldots, n$
  $$
  \dx{y}=(x+1)(x+2)(x+3)\ldots(x+n)\left[\frac{1}{x+1}+\frac{1}{x+1}+\cdots+\frac{1}{x+1}\right]\dx{x}$$
  Can you find a way to extend this so that it is also true for $x=1, 2,
  3, \ldots, n$
  
  \noindent{}\hint{This should look familiar. Did you do problem \# in Section \#}
\end{myproblem}
\end{ProblemSection}

 
   
\section{General Logarithms and Exponentials-tmp}
\label{sec:gener-logar-expon}
So far we've focused on the \emph{natural} exponential and logarithm
functions. But if the natural exponential is just this screwy number
$e$ raised to $x$ then isn't say, $2^x$ also an exponential?

Of course it is. It is the ``exponential with base $2$.''

So is $3^x$, It is the ``exponential with base $3$.'' Just as $5^x$ is
the ``exponential with base $5$'' and $32543^x$ is the ``exponential
with base $32543$.'' In fact for any $a>0$ then $a^x$ is \emph{an}
exponential.
\begin{ProficiencyDrill}
  Why do we impose the condition, $a>0?$
\end{ProficiencyDrill}
And  for every function of the form $a^x$ there is a corresponding 
inverse function called the ``log base $a$,'' and denoted $\log_a(x)$.

The following problem establishes that all logarithm functions change
multiplication into addition and exponentiation into multiplication as
we stated (but didn't prove) earlier.
\begin{embeddedproblem}{}
  Let $a>0,$ $x>0,$ $y>0$ and  let $c$ be any real number.
  \comment{Calculus is not needed for this problem.}
  \begin{enumerate}[label={\bf{}(\alph*)}]
  \item Show that $\log_a  xy=\log_a x+\log_a y.$
    \hint{Consider
      $a^{\log_a x+\log_a y}$  and use the properties of exponents to show that
      this equals $xy.$}
  \item Show that $\log_a x^c =c \log_a x.$
    \hint{Use the properties of exponents to show that both
      $a^{c\log_a(x)}$ and $a^{\log_a(x^c)}$ are equal to $x^c.$}
  \end{enumerate}
\end{embeddedproblem}

What makes the \term{natural exponential}, ``natural'' is that the choice
of $e$ as the base makes it is the only exponential which is its own
derivative: $\dfdx{(e^x)}{x}=e^x$.

But if there are so many different exponential and logarithm functions
don't we need differentiation rules for all of them? And what are the
derivatives of the other exponentials and logarithms?

\subsubsection*{\underline{General Exponentials}}
It is not as hard to differentiate $y=a^x$ as you might think. Since
we already know how to differentiate the \term{natural exponential}
all we have to do is re-express $y=a^x$ in terms of $e^x$. This is
usually called "Here's how.

We know that $a=e^{\ln(a)}$ so we make that substitution to get
$$
y=a^x=\left(e^{\ln(a)}\right)^x=e^{x\ln(a)}
$$
You can think of this step as making $a^x$ ``harder on the eyes'' if
you like. We're just making a substitution.

Now we differentiate:
\begin{align*}
  a^x&=e^{x\ln(a)}\\
  \dx(a^x) &= e^{x\ln(a)}\dx(x\ln(a)).\\
  \intertext{Since $a$ is a constant, $\ln(a)$ is also a constant so}
  \dx(a^x) &= e^{x\ln(a)}\ln(a)\dx{x}\\
  \intertext{and since $a^x=e^{x\ln(a)}$ we have}
    \dx(a^x) &= a^x\ln(a)\dx{x}\\
\end{align*}

\begin{ProficiencyDrill}
  Find $\dx{y}$ and $\dfdx{y}{x}$ for each of the following:
  \begin{enumerate}[label={\bf{}(\alph*)}]
    \begin{multicols}{3}
    \item $y=2^x$
    \item $y=3^x\cdot3^x$
    \item $y=x3^x$
    \item $y=3\cdot7^x+7\cdot3^x$
    \item $y=e^{5^x}$
    \item $y=\ln\left(2^x\right)$
    \item $y=2^x\cdot4^x\cdot8^{x^2}$
    \end{multicols}
  \end{enumerate}
\end{ProficiencyDrill}

\subsubsection*{\underline{General Logarithms}}
Just as we did for the general exponential $a^x$ we compute the
derivative of the logarithm by changing it to a natural logarithm
($\ln(e)$).

To do this we start with the exponential. If $y=\log_a(x)$ then
\begin{align*}
  a^y&=x,\  a>0.\\
  \intertext{Taking the \term{\underline{natural} logarithm} and using
  Property 4 from out list of properties  we have,}
  y\ln(a)&=\ln(x), \text{ or}\\
  y&=\frac{\ln(x)}{\ln(a)}\\
  \intertext{Since $y=\log_a(x)$ we have successfully converted the general
  logarithm to a base-$a$ logarithm:}
  \log_a(x)&=\frac{\ln(x)}{\ln(a)}.\\
\intertext{
An easy mnemonic for remembering this formula is to simply notice that
  on both sides of the equation the $a$ appears lower than the $x$.}
  \intertext{
 Since $\ln(a)$ is a
 constant (why?) we have}
  \dx{(\log_a(x))}&=\frac{1}{\ln(a)}\cdot\dx\left(\ln(x)\right)\\
     &=\frac{1}{\ln(a)}\cdot\frac1x\dx{x}.\\
  \intertext{So we have the differentiation formula:}
  \dfdx{(\log_a(x))}{x}&=\frac{1}{x\ln(a)}.
\end{align*}

\subsubsection*{\underline{Common Logarithms}}
It is a well known phenomenon in mathematics that the first solution
of a problem is often the most difficult, convoluted, and hard to
understand solution possible. This is what happened to Napier. His
original scheme, although it did allow him to convert multiplication
to addition, was more complicated than necessary.  Shortly after
Napier published his work in $1614$, Henry Briggs ($1561 - 1630$),
professor of geometry at Gresham College in London, visited Napier and
suggested that a few simple changes to his logarithm table would make
it more practical. One change Briggs suggested was to use a table of
base-$10$ logarithms, since our number system is in base-$10$. Napier
had had the same idea himself, but had been unable to pursue it ``on
account of ill-health.'' As a result Briggs took on the task of
computing tables of base-$10$ logarithms, collaborating with Napier
until the latter's death two years later. The result of their
collaboration known as the Briggsian or \term{Common
  Logarithms}. Common Logarithms were an important computational aid for
slightly over $300$ years until the invention of modern computational
technologies. 

As the mathematical historian Howard Eves put it, the invention of
logarithms literally ``doubled the life of the astronomer,'' because
they so drastically reduced the time spent doing arithmetic.

Using Briggs's tables if a scientist  needed to
multiply, for example, $21.0234\times107.889766$ (s)he would first look up the
logarithm of $21.0234$ and $107.889766$ in his table, add the two
logarithms together, and then turn to his table again to see what the
sum was the logarithm of (this was called the 'anti'-logarithm). That
would be the product. 

Using his tables taking roots is also straight forward because
logarithms not only turned multiplication into addition they turn
exponentiation into multiplication. For example, to find compute
$\sqrt[2]{4}$ we would look up its logarithm, divide that by
three, and then look up the 'anti'-logarithm of the result.  Think
about it, would you rather find the cube root of $31$ by hand (do you
even know how to do that?) or divide its logarithm by $3$, and then
look up the root in a table?

Most modern scientific calculuators have both the natural common
logarithms built in. Naturally the ``\term{ln}'' button on a calulator
compute natural logarithms while the ``\term{log}'' computes common
logarithms.
\begin{ProficiencyDrill}
  \begin{enumerate}[label={\bf{}(\alph*)}]
  \item Use whatever technology you prefer to compute the natural
    logarithms of the numbers $21.2343$ and $5689.121343$ and their
    product. Confirm that
    $\ln(21.2343) \ln(5689.121343) =
    \ln(21.2343\times5689.121343)$. \comment{If you don't like these
      numbers use others. We just picked these at random.}
  \item Do the same using common logarithms.
  \item Do the same using $\log_5(x)$. \comment{Or any other
      base. Again we  just picked $5$ at random.}
  \end{enumerate}
\end{ProficiencyDrill}


\begin{embeddedproblem}
  How many base-$10$ digits long do you suppose the number
  $2^{1234567890}$ is? You cannot solve this by punching
  $2^{1234567890}$ into a calculator and counting the digits. As
  powerful and convenient as our modern technology is there are still
  problems, like this one, that require a deeper understanding. A
  calculator is a useful tool, but the human operating the
  calculator needs to understand what to calculate. Blind
  calculation won't work is pointless.

  Suppose we let $n = 2^{1234567890}$ and take the base $10$ logarithm
  (common logarithm) of both sides. This gives:
  $$
  \log_{10}(n)= \log_{10}(2^{1234567890})=1234567890\cdot
  \log_{10}(2)\approx371641966.574.
  $$
  Thus
  $$
  2^{1234567890}=n=10^{371641966.574}=10^{0.574}\times10^{371641966}
  $$
  So our number is $371641966$ base-$10$ digits long. Do you see why?

  This feels like a problem made up by math professors just to torment
  our students, but it is not. Computer programmers routinely have to
  allocate space in memory to hold information. If the information
  being held happens to be the number  $2^{1234567890}$ they would
  need to know how much space to allocate. That computation is
  essentially the one we just did.

  Of course, since computer arithmetic is binary (base-$2$) they would
  need to know how many binary digits (bits) to allocate. Obviously that
  would be $1234567890$ bits. Do you see why?

  But what about $10^{1234567890}?$ How many bits are needed to store
  this number?
\end{embeddedproblem}
  
\begin{embeddedproblem}
  For each of the following numbers:
  \begin{enumerate}[label={\bf{}(\roman*)}]
      \begin{multicols}{3}
      \item $3^{1234567890}$
      \item $9^{1234567890}$
      \item $\left(\frac13\right)^{1234567890}$
      \item $\left(\frac19\right)^{1234567890}$
      \end{multicols}
    \end{enumerate}
  \begin{enumerate}[label={\bf{}(\alph*)}]
    \item Find the number of binary digits needed to store the number.
    \item Find the number of base-$8 $ digits needed to store the number.
    \item Find the number of base-$10$ digits needed to store the number.
    \item Find the number of base-$100$ digits needed to store the number.
    \item Find the number of base-$9  $ digits needed to store the number.
  \end{enumerate}
\end{embeddedproblem}

% \begin{enumerate}[label={\bf{}(\alph*)}]
%     \item in base-$10$
%     \item in base-$3$
%     \item in base-$2$
%     \item in base-$7$
%     \end{enumerate}
%   \end{embeddedproblem}

\subsubsection*{\underline{Napierian Logarithms}}
The discussion above begs the question: What base did Napier use when
he invented logarithms?

Napier defined his original logarithms as follows.  Consider a point
$C$ moving along line segment $AB$ and a point $F$ moving along an
infinite ray
$DE$ as seen below.\\
\centerline{\includegraphics*[height=1in,width=3in]{../Figures/NapierLog1}}
Assume that they start at the same speed and that $F$ moves at a
constant speed, but that $C$ moves with a speed equal to the distance
from $C$ to $B$.  If we let $CB=x$ and $DF=y$, then we define the
\term{Naperian Logarithm} as
$$
y = \text{Nap log}(x).
$$

Since Calculus had't been invented yet Napier used Trigonometry to
develop his \pubtitle{Table of Logarithms}. He took the length of
$AB$ to be $10^7$, since the best sine tables of the day were to
accurate to 
seven decimal places.



\begin{embeddedproblem}
  It is not altogether clear from the description above that  the
  Napierian logarithm actually is a logarithm, let alone what its  base
  is. This question explores both of these questions.

  \label{pic:NapierianLog}
  \begin{enumerate}[label={\bf{}(\alph*)}]
  \item  Show that $\dfdx{y}{t}=10^7$, $y(0)=0$.
  \item Solve the IVP: $\dfdx{x}{t}=-x$,  $x(0)=10^7$,
  \item Use the information in parts (a) and (b) to show that
    $$
    \text{Nap log}(x)=y(x)=10^7\log_{\frac{1}{e}}\left(\frac{x}{10^7}\right).
    $$
    \comment{In the formula above we have suppressed the variable $t$.}
  \end{enumerate}

\end{embeddedproblem}




\subsubsection*{\underline{Logarithmic Differentiation}}
The ability to turn multiplication into addition and exponentiation into
multiplication can make some differentiations much easier to do.

\begin{myexample}
  For example, suppose we need to differentiate
$$
y=\frac{\sqrt{x^3+2x}\cdot{}\tan(x)}{(x^5-7)^4}.
$$

Certainly we can do this using  our rules, but it will clearly be
tedious.  But nice things happen if we take the natural logarithm of both sides before we
differentiate.
\begin{align*}
  \ln(y)&=\ln\left(\frac{\sqrt{x^3+2x}\cdot{}\tan(x)}{(x^5-7)^4} \right)\\
%        &=\ln\left(\frac{x^3+2x)^{1/2}\cdot{}\tan(x)}{(x^5-7)^4}\right).
\intertext{Since the natural  logarithm changes division   into subtraction we
          can rewrite the right side as.}
         &=\ln\left[\sqrt{x^3+2x}\cdot\tan(x)\right] -
             \ln\left(\left(x^5-7\right)^4\right).
           \intertext{And since the  natural  logarithm changes
             multiplication    into addition we           can
           re-express  the right side again. This time as:}
         &=\ln\left(\sqrt{x^3+2x}\right)+ \ln\left(\tan(x)\right) -
           \ln\left(\left(x^5-7\right)^4\right).
           \intertext{Re-expressing the square root as an exponent
           gives}
                   &=\ln\left(x^3+2x\right)^{1/2}+ \ln\left(\tan(x)\right) -
                     \ln\left(\left(x^5-7\right)^4\right).
                     \intertext{and we can now bring the exponent down
                     in front. We're not using the Power Rule
                     here. We're using the fourth property of
                     logarithms in our table.} 
     \ln(y)    &=\frac{1}{2}  \ln(x^3+2x)+\ln(\tan(x)) -4 \ln(x^5-7)).
 \end{align*}
Notice that all we've done so far is re-express our function using the
properties of logarithms. We have not yet started differentiating.

But now differentiating this is relatively easy.
\begin{align*}
  \dx{(\ln(y))}&=
                 \frac{1}{2(x^3+2x)}\dx{(x^3+2x)}+\frac{1}{\tan(x)}\dx{(\tan{x})}-\frac{4}{x^5-7}\dx{(x^5-7)}\\
  \frac{\dx{y}}{y}&=
                    \frac{1}{2(x^3+2x)}(3x^2+2)\dx{x}+\frac{1}{\tan(x)}\sec^2(x)\dx{x}-\frac{4}{x^5-7}(5x^4)\dx{x}\\
\end{align*}
If we multiply by $y$ we get
\begin{align*}
  \dx{y}&=
          y\left[\frac{(3x^2+2)}{2(x^3+2x)}\dx{x}+\frac{\sec^2(x)}{\tan(x)}\dx{x}-4\left(\frac{5x^4}{x^5-7}\right)\dx{x}\right]\\
        &=        \left[\frac{\sqrt{x^3+2x}\cdot{}\tan(x)}{(x^5-7)^4}\right]\cdot\left[\frac{(3x^2+2)}{2(x^3+2x)}\dx{x}+\frac{\sec^2(x)}{\tan(x)}\dx{x}-4\left(\frac{5x^4}{x^5-7}\right)\dx{x}\right]\\
\end{align*}
\end{myexample}

\begin{ProficiencyDrill}
  This looks like a lot of work when it's laid out on the page like this
but it really isn't. With practice you can differentiate an expression
like 
$
\frac{\sqrt{x^3+2x}\cdot{}\tan(x)}{(x^5-7)^4}.
$ in your head  as fast as you can write it down. Really.

Even if this is not true,   would
you rather compute this derivative using the quotient rule and
product rule?

  Try it and see.
\end{ProficiencyDrill}

This technique of taking a natural logarithm before you
differentiate is called logarithmic differentiation, and sometimes it
is an alternative, especially if you have a lot of product rules
and/or quotient rules to deal with.  You could actually do it all the
time, but sometimes, it is just not worth it.  For example, you could
use it for differentiating $y=x^2+1.$ 
\begin{align*}
  y&=x^2+1\\
  \ln(y)&=\ln(x^2+1)\\
  \frac{1}{y} \dx{y}&=\frac{1}{x^2+1} 2x\dx{x}\\
  \dx{y}&=y\left(\frac{2x}{x^2+1}\right)\dx{x}\\
   &=(x^2+1)\left(\frac{2x}{x^2+1}\right)\dx{x}\\
   &=2x \dx{x}.
\end{align*}
It gets you the right answer, but wouldn't it have been easier to just
the Constant Multiple Rule, and then the Power Rule?  So use your
judgement.  Use a logarithm if it makes your job easier.

But there are some situations where logarithmic
differentiation is your only option.  Consider something like $y=x^x.$ 
Unfortunately, this is not the power rule like $\dx{x}{x^2}$ nor is it an
exponential like $\dx{x}{2^x }$.  It is some sort of strange
combination of both, so our existing rules don't directly apply.
However, this does not say that it can't be differentiated.  This is
the sort of situation where logarithmic differentiation is just the
technique we need. 

\begin{embeddedproblem}
  \begin{enumerate}[label={\bf{}(\alph*)}]{}
  \item Apply logarithmic differentiation to $y=x^x$ to obtain
    $$
    \dx{(x^x) }=x^x (1+\ln(x) )\dx{x}.
    $$
  \item  Given $y=(x^2+1)^{\sin(x)},$ compute $\dx{y}.$
  \end{enumerate}
\end{embeddedproblem}

% Since we defined the exponential as the solution of the differential
% equation
% $$
% \dfdx{y}{x}=y
% $$
% it follows that the exponential is its own derivative. Similarly since
% the logarithm is a solution of the differential equation
% $$
% \dfdx{y}{x}=\frac1x
% $$
% it follows that $\dfdx{\ln(x)}{x} = \frac1x.$


% \subsection*{\underline{General Logarithms}}
% The function which ``undoes'' the base-$2$ exponential is called the
% base-$2$ logarithm and is denoted: $\log_2(x)$. In general the inverse
% of $a^x$, $a>0$ is $\log_a(x)$. The next question for us is: What is
% the derivative of $\log_a(x)?$

% As before this just a matter of changing the base to $e$, but it is
% slightly more complicated to change the base of a logarithm so we
% address that first.

% Suppose we want to re-express $\log_5(x)$ as a \term{natural
%   logarithm} (logarithm base $e$). Let $y=\log_5(x).$ Then clearly
% \begin{align*}
%   5^y&=x \text{ This is clear since $\log_5(5^y)=y$.}\\
%   \intertext{Taking the \term{natural logarithm} of both sides we
%   have}
%   \ln(5^y)&=\ln(x)\\
%   y\ln(5)&=\ln(x) \text{ so that}\\
%   y&=\frac{\ln(x)}{\ln(5)}.
% \end{align*}

% Since $y=\log_5(x)$ we have 
% $$
% y=\log_5(x) =\frac{\ln(x)}{\ln(5)}, 
% $$ 
% and we have expressed $\log_5(x)$ as a natural logarithm.

% Since we know how to differentiat the natural logarithm the rest is
% straightforward:
% $$
% \dx\left(\log_5(x)\right) =\frac{1}{x\ln(5)}\dx{x}.
% $$


%%% Local Variables: 
%%% mode: latex
%%% outline-minor-mode: t
%%% eval:(flyspell-mode)
%%% TeX-master: "DiffCalc"
%%% End: 
